% 5.1 — Le travail offensif
\section{Le travail offensif}
\label{sec:offensif}

\subsection{Trois façons de chercher}
Les failles ont été cherchées comme dans un programme de \textit{bug bounty}, avec trois
méthodes complémentaires :
\begin{description}[style=nextline,leftmargin=0.6cm,font=\sffamily\bfseries]
  \item[La lecture de code] Lire un module en se demandant « que se passe-t-il si cette
    hypothèse est fausse ? » --- par exemple, si le fichier que la tâche élevée exécute est
    inscriptible par l'utilisateur.
  \item[Les sondes] Fabriquer des commandes piégées et les faire juger dans un projet jetable :
    une substitution qui lit un fichier SQL, un outil qui supprime son entrée sans que son
    nom l'annonce, une valeur qui commence par un tiret.
  \item[La fouille des registres] Analyser les reçus réellement produits par les agents pour
    voir ce qu'ils exécutent et comment \demo{} les a jugés. L'analyse menée pour C3 a porté
    sur \nbregistres{} registres et \nbrecus{} reçus, et a montré le poids réel des lignes
    PowerShell jugées par les anciennes règles.
\end{description}
Chaque faille confirmée suit la même règle que le reste du projet : un test qui \emph{échoue}
sur le code vulnérable est écrit avant la correction.

\subsection{L'audit de reprise}
Au début de la période documentée, un audit de la version 1.7.0 a confirmé 18 défauts, plus un
contournement du garde du shell par auto-exclusion. Ils ont d'abord été \emph{figés} par des
tests rouges --- la suite passait alors avec 1856 tests réussis, aucun ne décrivant ces
comportements --- puis corrigés en trois tours. Ils se rangent en trois niveaux :
\begin{itemize}
  \item \textbf{l'outil revendique une récupération qu'il n'a pas prise} : une modification de
    fichier hors du projet déclarée REVERSIBLE, un lien symbolique restauré comme un fichier
    ordinaire, un faux répertoire \fichier{.demo\_cli} ignoré par son \emph{nom}, un
    \cmd{cp -f} sur un fichier existant non capturé ;
  \item \textbf{le garde est absent, contournable ou muet} : un préfiltre du shell plus étroit
    que le classifieur, un appel au garde sans délai maximal, une panne qui ne laisse aucune
    trace, un interrupteur d'arrêt non relu dans un shell déjà ouvert, une auto-exclusion
    décidée sur une sous-chaîne au lieu du mot de commande ;
  \item \textbf{du texte jugé comme une commande} : un message de commit ou un corps de
    \textit{heredoc} pris pour des commandes, une aide (\cmd{--help}) prise pour une
    modification.
\end{itemize}

\subsection{La batterie red-team}
Les notes de red-team tenues pendant le stage classent chaque constat en trois catégories, qui
découlent directement du modèle de menace (section~\ref{sec:menace}) :
\begin{itemize}
  \item \textbf{défaut dans le périmètre} (\textit{in-scope flaw}) : l'agent commet une erreur
    honnête et le garde la manque --- c'est un défaut à corriger ;
  \item \textbf{frontière} (\textit{boundary}) : une limite documentée de ce que le \textit{hook}
    peut voir ;
  \item \textbf{évasion} (\textit{evasion}) : un obscurcissement délibéré, hors du périmètre par
    conception.
\end{itemize}

{\small
\begin{longtable}{@{}L{0.9cm}L{6.3cm}L{2.4cm}L{5.0cm}@{}}
\caption{La batterie red-team \#001--\#010 et son issue}\label{tab:redteam}\\
\toprule
\textbf{\#} & \textbf{Constat} & \textbf{Catégorie} & \textbf{Issue} \\
\midrule
\endfirsthead
\toprule
\textbf{\#} & \textbf{Constat} & \textbf{Catégorie} & \textbf{Issue} \\
\midrule
\endhead
\bottomrule
\endfoot
001 & Les commandes tapées dans le mode \cmd{!} de Claude Code n'atteignaient pas le
  \textit{hook} (elles s'exécutent dans le shell de l'utilisateur) & Frontière &
  Couvertes par le garde du shell (piège \cmd{DEBUG} injecté par \cmd{BASH\_ENV}) ; les commandes
  tapées dans d'autres shells restent hors du filet Linux (V14) \\
002 & Des commandes obscurcies (Bashfuscator) cachent les mots destructeurs à un classifieur
  textuel & Évasion & Hors périmètre de la lecture statique ; le filet d'exécution voit les vrais
  appels système \\
003 & \cmd{source fichier.txt} exécute des commandes que le \textit{hook} ne voit pas &
  Frontière & Traité comme du code non lu : filet, sinon point de contrôle, sinon ALLOW avec une
  note honnête \\
004 & \textbf{Faux REVERSIBLE} : une suppression avec une expansion non résolue ne capturait que
  la partie résolue et déclarait toute la ligne réversible & Défaut & Corrigé : une capture
  partielle est refusée et la ligne escalade \\
005 & Redirections et verbes qui écrasent (\cmd{> fichier}, \cmd{tee}, \cmd{cp} sur un fichier)
  détruisaient sans déclencher de règle & Défaut & Corrigé : lus par le moteur ; créer un
  nouveau fichier reste autorisé \\
006 & Un préfixe de variable d'environnement (\cmd{X=1 rm app.db}) contournait la règle
  \cmd{rm} ancrée & Défaut & Corrigé \\
007 & Trous de couverture : \cmd{mkfs}, \cmd{git checkout -- .}, \cmd{git restore .},
  \cmd{git stash clear}, \cmd{Remove-Item} sans \cmd{-Force} & Défaut & Corrigé \\
008 & Mesure sur 88 incidents réels : couverture de 75\,\% par enregistrement et 44\,\% des
  commandes uniques & Mesure & Portée à 90\,\% et 81\,\% après les corrections \\
009 & La destruction validée est rare dans les données réelles ; deux trous : les écritures SaaS
  qu'aucune copie locale n'annule, et \cmd{git stash pop} jugé sûr & Frontière + défaut &
  Écritures SaaS escaladées comme surfaces non récupérables ; \cmd{git stash pop} corrigé \\
010 & Essais en direct du mode \cmd{!} : l'enveloppe \cmd{eval '<cmd>'} prise pour une
  exécution distante, une redirection vers un nouveau fichier prise pour une destruction, chaque
  commande capturée deux fois & Défaut $\times$ 3 & Corrigés \\
\end{longtable}
}

Au fil des parties, d'autres défauts ont été trouvés et corrigés ; quelques-uns méritent une
ligne :
\begin{itemize}
  \item \textbf{Garde Windows éteint en silence} : \cmd{Out-File} de PowerShell écrit en UTF-8 avec
    BOM, la configuration ne se lisait plus et le garde laissait tout passer. Une configuration
    illisible escalade désormais, avec un message qui dit comment la réparer.
  \item \textbf{Montage WinFsp} : confusion entre suppression de fichier et de répertoire,
    renommage sur un fichier existant non capturé, descripteurs non fermés, restauration écrite
    au-delà du montage, verrous de plages d'octets à travers le montage, point de montage périmé
    après un redémarrage.
  \item \textbf{Intégrité de \cmd{undo}} : une suppression récursive restaurée comme un répertoire
    vide (Linux), une restauration qui échouait sur un parent manquant (Windows), des lignes
    déchirées dans l'index de récupération en écriture concurrente (corrigé par un verrou
    inter-processus).
  \item \textbf{Points de contrôle} : un calcul de taille qui ignorait \fichier{.git} et
    permettait de dépasser la limite de 256\,Mo ; des copies complètes de bases de plusieurs Go
    qui figeaient la machine (liste d'exclusion partagée et limites strictes).
  \item \textbf{Installation} : un filtre \cmd{Bash|PowerShell} comparé comme une chaîne entière
    déclenchait le \textit{hook} deux fois par commande ; la désinstallation retirait les
    \textit{hooks} d'autres outils.
  \item \textbf{Faux positifs} : un mot comme « drop » dans un message de commit déclenchait les
    règles SQL ; un aperçu de \cmd{cd x \&\& rm -rf a b} montrait le mauvais répertoire.
\end{itemize}

\subsection{Les vulnérabilités du stage}
Le journal des vulnérabilités (V1--V22) distingue ce qui est corrigé, ce qui a été vérifié et
n'est pas une faille, et ce qui reste ouvert.

{\small
\begin{longtable}{@{}L{1.5cm}L{6.3cm}L{1.6cm}L{1.9cm}L{3.3cm}@{}}
\caption{Vulnérabilités corrigées}\label{tab:v-corrigees}\\
\toprule
\textbf{Id} & \textbf{Description} & \textbf{Gravité} & \textbf{Trouvée par} & \textbf{Correction} \\
\midrule
\endfirsthead
\toprule
\textbf{Id} & \textbf{Description} & \textbf{Gravité} & \textbf{Trouvée par} & \textbf{Correction} \\
\midrule
\endhead
\bottomrule
\endfoot
V1 & La tâche Windows d'ouverture de session (niveau « le plus élevé ») exécutait un script de
  \cmd{\%LOCALAPPDATA\%} modifiable par l'utilisateur, donc par l'agent : code exécuté avec
  élévation, ou garde désactivé par un \cmd{exit /b 0}. & Haute & Lecture de code, vérifié
  sous Windows & W-SEC (\cmd{db262df}) \\
V2 & Même sans ce script, le code de la tâche était inscriptible : installation éditable et
  environnement virtuel utilisateur. & Haute & Inspection & W-SEC \\
V3 & La tâche trouvait \cmd{demo\_cli} par le \cmd{PATH} de l'utilisateur, modifiable. & Haute &
  Lecture de code & W-SEC \\
V4 & La tâche élevée écrivait son journal dans un dossier inscriptible par l'utilisateur
  (redirection par un lien). & Moyenne & Revue de conception & W-SEC \\
V17 & Points de contrôle activés : une ligne illisible avec un opérande \emph{hors} du projet
  recevait une copie du projet et REVERSIBLE --- une promesse non tenue. & Haute (option) &
  Lecture, reproduit & V17 (\cmd{99831a9}) \\
V18 (1) & Les filets et le point de contrôle couvrent le projet seulement, mais la raison disait
  « chaque fichier détruit est sauvegardé ». & Moyenne (honnêteté) & Lecture de code & M2.3a \\
V19 & SQL non lu autorisé : \cmd{psql -c "\$(cat q.sql)"} avec \cmd{DROP TABLE users;} $\rightarrow$
  ALLOW, rien capturé. & Haute & Sondes & M2.3a (\cmd{3c94414}) \\
V20 & Outils de fichiers destructeurs non vus : \cmd{gzip a.txt}, \cmd{xz}, \cmd{bzip2},
  \cmd{sort -o}, \cmd{iconv -o}, \cmd{patch}, \cmd{awk -i inplace}, \cmd{7z \ldots{} -sdel}
  $\rightarrow$ ALLOW. & Haute & Sondes & M2.3a \\
V21 & Un exécuteur de code auquel on passait un fichier hors du projet (\cmd{python x.py
  /etc/hostname}) perdait le filet, disponible pourtant pour \cmd{python x.py}, et devenait ALLOW. La note affichée restait honnête. &
  Moyenne & Sondes & M2.3b (\cmd{ad7ab16}) \\
V22 & La suppression d'un lien symbolique du projet était capturée, mais \cmd{undo} le
  restaurait comme un \emph{fichier ordinaire} : la restauration utilisait l'entrée de groupe qui
  enveloppe une cible unique, au lieu de l'entrée interne qui garde le lien. Pas de perte de
  données. & Moyenne & Revue
  red-team (R3) & V22 (\cmd{7d2a2aa}) \\
V12 & PowerShell non analysé : \cmd{iex "Remove-Item a.txt"}, \cmd{Write-Output "\$(Remove-Item
  a.txt)"} et \cmd{"a.txt" | Remove-Item} autorisés, \cmd{Copy-Item a b} écrasait en silence,
  \cmd{Set-Content a.txt hi} escaladait inutilement. & Haute & Fouille des registres & C3 \\
V18 (2) & Une destruction hors du projet par du code non lu n'était ni sauvegardée ni
  enregistrée par le filet. & Moyenne & Revue de code (M2.3a) & V18 (\commitVdixhuit) : le filet
  Linux enregistre chaque fichier détruit hors du projet \\
--- & \cmd{source ./x.sh} vu comme sans effet ; code téléchargé exécuté par
  \cmd{bash <(curl \ldots)} ou \cmd{aws s3 cp \ldots{} - | bash} autorisé. & --- & Sondes & M2.3a \\
--- & \cmd{undo} impossible sur un fichier passé en \cmd{chmod 444}/\cmd{000} ; plantage du
  moteur sur \cmd{git stash pop \&\& git stash drop} ; \fichier{\textasciitilde/.bashrc}, qui
  active le garde du shell, modifiable par l'agent. & --- & Sondes & M2.2c, M2.2d \\
\end{longtable}
}

Deux points ont été vérifiés et ne sont \emph{pas} des failles : un utilisateur normal ne peut
pas modifier la tâche élevée elle-même (V5, « Access is denied ») ; \cmd{Move-Item} et
\cmd{Rename-Item} refusent d'écraser une destination existante (vérifié sur PowerShell 5.1) ---
seul \cmd{Copy-Item} écrase, ce qui relève de V12.

Un point est un \emph{risque préventif} plutôt qu'une faille active. \textbf{V16} (gravité élevée
s'il s'était réalisé) : en retirant le pouvoir de décision aux anciennes règles, M2.3 risquait de
perdre le refus du point de contrôle pour une commande visant un fichier hors du projet
(\cmd{/etc}, \cmd{\textasciitilde/.ssh}) ; elle aurait reçu une copie du seul projet, présentée à
tort comme réversible. M2.3a et M2.3b ont conservé ce refus, épinglé par les tests de V17 et de
V21.

{\small
\begin{longtable}{@{}L{1.8cm}L{8.2cm}L{4.6cm}@{}}
\caption{Vulnérabilités et constats ouverts}\label{tab:v-ouvertes}\\
\toprule
\textbf{Id} & \textbf{Description} & \textbf{Plan} \\
\midrule
\endfirsthead
\toprule
\textbf{Id} & \textbf{Description} & \textbf{Plan} \\
\midrule
\endhead
\bottomrule
\endfoot
V14 & Les commandes tapées par la personne et les processus qui survivent à la commande sont
  hors du filet Linux. & Plus tard \\
V7--V9 & Montage Windows : ne refuse pas lui-même les fichiers de \demo{} (les portes le
  font), lit les fichiers entiers en mémoire, santé réduite à « processus vivant ». & Partie
  « parité Windows » \\
V10--V11 & Proxy réseau : repli silencieux si son outil manque ; périmètre d'exfiltration
  (un \cmd{curl -X POST -d @-} est escaladé, d'autres outils non modélisés). & Partie réseau \\
V15 & Codex ne peut pas exécuter une commande réécrite sans l'autoriser : pas de filet sur
  Codex. & Portage Codex \\
V6, V13 & Reçus : rares lignes déchirées quand Windows et WSL écrivent le même registre ; pas de
  rotation du registre ni de la réserve. & Maintenance \\
\end{longtable}
}
